\chapter{跨接收机失配诊断与 RCPA-T 校准}
\label{ch:ch4}

本章对应 Paper 2，协议为 \textbf{block-disjoint cross-receiver}，与第\ref{ch:ch3}章 cross-day 不同。

\section{跨接收机问题定义}

源接收机 $R_s$ 训练、目标接收机 $R_t$ 测试。Source-only 闭集约 19--21\%，接近 chance。目标：在少量目标域标签下恢复可用识别性能。

\section{接收机诱导的 OOB 特征纠缠}

跨接收机时，OOB 同时承载设备与接收机信息。诊断证据（图~\ref{fig:ch4_diag}）：

\begin{itemize}
  \item OOB 接收机探针准确率 \textbf{72.7\%}；
  \item 跨 RX fused 嵌入距离比 collapse（约 0.22）；
  \item 预测分布高度集中（top-1 mass 约 95.8\%）。
\end{itemize}

\begin{figure}[htbp]
  \centering
  \includegraphics[width=0.95\linewidth]{fig1_diagnosis_summary.pdf}
  \caption{跨接收机诊断：谱偏置、probe、距离比与预测 collapse}
  \label{fig:ch4_diag}
\end{figure}

\section{RCPA-T 方法}

冻结第\ref{ch:ch3}章主干与源分类器；用每设备 $K$ 个 labeled target windows 估计目标接收机类原型 $\mu_c^{(t)}$；推理时最近原型分类。Support/query 在块级不重叠。

\section{主实验结果}

表~\ref{tab:ch4_baseline}：source-only near chance。表~\ref{tab:ch4_rcpa} 与图~\ref{fig:ch4_shot}：RCPA-T pooled $K{=}5/10/20$ 为 \textbf{58.3\% / 69.4\% / 75.0\%}。

\begin{table}[htbp]
  \centering
  \caption{跨接收机 source-only 基线（文件级准确率，\%）}
  \label{tab:ch4_baseline}
  \begin{tabular}{lcc}
    \toprule
    方向 & Source classifier & RCPA-S \\
    \midrule
    RX1$\rightarrow$RX2 & $19.4\pm3.4$ & $15.3\pm7.3$ \\
    RX2$\rightarrow$RX1 & $20.8\pm8.8$ & $13.0\pm3.1$ \\
    \bottomrule
  \end{tabular}
\end{table}

\begin{table}[htbp]
  \centering
  \caption{RCPA-T shot curve（pooled，文件级准确率，\%）}
  \label{tab:ch4_rcpa}
  \begin{tabular}{cccc}
    \toprule
    $K$ & RX1$\rightarrow$RX2 & RX2$\rightarrow$RX1 & Pooled \\
    \midrule
    5 & $57.4\pm9.4$ & $59.3\pm8.5$ & \textbf{58.3$\pm$9.0} \\
    10 & $67.6\pm10.0$ & $71.3\pm9.1$ & \textbf{69.4$\pm$9.7} \\
    20 & $77.3\pm7.1$ & $72.7\pm8.1$ & \textbf{75.0$\pm$8.0} \\
    \bottomrule
  \end{tabular}
\end{table}


\begin{figure}[htbp]
  \centering
  \includegraphics[width=0.85\linewidth]{fig2_rcpa_shotcurve.pdf}
  \caption{RCPA-T K-shot 曲线}
  \label{fig:ch4_shot}
\end{figure}

同协议 K=5 时 linear probe 59.0\% 略高于 RCPA-T 58.3\%（约 0.7 pp，在 seed 方差内）。无标签对齐 baseline 约 20--25\%，无效。

\section{局限}

需 labeled target windows；非 source-free；仅 RX1/RX2；未覆盖第\ref{ch:ch5}章 EM 与 open-set。

\section{本章小结}

本章诊断 OOB 纠缠并提出 RCPA-T，在少样本目标校准下显著优于 source-only，引出复杂电磁与开放集场景的第\ref{ch:ch5}章。
